\scanpage{85}
(v) On n'est pas dans le ``cas \deleted{\unsure{aberrant}} \inserted{impropre}'' suivant \inserted{$(\alpha=0,\beta=-2)$}
\[ \left\{\begin{aligned}&\deleted{\unsure{p\ne2}}\ \operatorname{car} k=p\ne2,\quad \sigma_0=\begin{pmatrix}-1&0&1\\0&-1&0\\0&0&1\end{pmatrix}\quad\text{i.e }\deleted{\sigma_0\text{ est}}\\
&\deleted{\text{la symétr}\unsure{\text{ie}}}\\
&\sigma_0(x,y)=(1-x,-y)\quad\text{i.e }\sigma_0=\text{symétrie p.r. au}\\
&\text{milieu du segment }(s_0,s_1)\end{aligned}\right. \tag*{(23)} \]

\uline{Dém} L'équivalence de (i), (ii), (ii') résulte
$\text{trivial}^{t}.$ de
\[ s_2-s_{-1}=\deleted{\unsure{\beta}}(1+\alpha)e_1+\beta e_2 \]
et du fait que $e_1=s_1-s_0$ et $e_2$ font une
base de $V$. L'équivalence de (ii) et (iii),
ou (ii') et (iii'), s'obtient en \deleted{\unsure{faisan|faisant}} notant
que pour $n=-1$ on obtient les conditions de
(ii) (iii), et qu'inversement (ii') (iii') s'obtiennent
de (ii) (iii) en appliquant $u^{n+1}$. On a
\[ \deleted{\unsure{\det\sigma_{0V}=}}\left\{\begin{aligned}\det u_V&=-\det\sigma_{0V}=1+\beta\\
\operatorname{Tr}\sigma_{0V}&=\beta\end{aligned}\right. \tag*{(24)} \]
ce qui prouve l'équivalence de (iv) \unsure{:|.} (iv') (iv'')
avec les autres conditions. Enfin, l'équivalence
avec (v) est évidente par la \deleted{\unsure{proposition|propos}}
formule (22).

\uline{Remarques sur le cas impropre} À isom unique
près, il y en a un et un seul pour \unsure{tt} $k$
% archive sheet number: 84
